Le plan d'action proposé détaille les modalités de conception, de déploiement et d'exploitation de ce questionnaire, dans une logique de gestion de projet étendue jusqu'à septembre 2027.

\subsection{Planification et structuration des tâches}

La réalisation du questionnaire s'inscrit dans une démarche progressive, structurée autour de plusieurs phases complémentaires :

\begin{enumerate}
    \item \textbf{Phase de cadrage stratégique}
    \begin{itemize}
        \item Sélection de trois \AssociationsRegionalesIHEDN{} pilotes afin de tester et valider le dispositif ;
        \item Réalisation d'une étude bibliographique approfondie (octobre 2025 - avril 2026) permettant de consolider le cadre conceptuel et méthodologique ;
        \item Obtention d'une validation institutionnelle formelle (\UnionIHEDN, Bureau, \acrshort{codir}), condition préalable à tout déploiement.
    \end{itemize}

    \item \textbf{Phase de conception du questionnaire}
    \begin{itemize}
        \item Définition de l'objectif général, des thématiques et de la problématique centrale ;
        \item Élaboration d'une trame structurée et d'une grille d'exploitation des données ;
        \item Conception du contenu détaillé et d'une maquette du questionnaire\footnote{cf. annexe \ref{annexe:maquette_questionnaire}} ;
        \item Choix des outils techniques et du support de diffusion.
    \end{itemize}

    \item \textbf{Phase d'identification des ressources}
    \begin{itemize}
        \item Identification des compétences nécessaires (experts en sciences sociales, spécialistes RGPD, partenaires académiques) ;
        \item Structuration des partenariats (\UnionIHEDN, IRSEM, institutions académiques) ;
        \item Estimation des ressources et planification détaillée.
    \end{itemize}

    \item \textbf{Phase de lancement et de déploiement}
    \begin{itemize}
        \item Mise en place d'un comité de pilotage ;
        \item Communication interne et externe sur les objectifs du projet ;
        \item Préparation et diffusion de la campagne par l'intermédiaire de courriers électroniques ;
        \item Administration du questionnaire auprès des membres.
    \end{itemize}

\clearpage

    \item \textbf{Phase d'exploitation et de valorisation}
    \begin{itemize}
        \item Traitement et analyse des données collectées ;
        \item Production d'indicateurs et de cartographies de compétences ;
        \item Rédaction d'un rapport de travail et capitalisation des résultats ;
        \item Restitution finale prévue le 15 septembre 2027.
    \end{itemize}
\end{enumerate}

Cette organisation suit un cycle complet de projet, de la conception à la valorisation, en cohérence avec l'objectif du comité : renforcer la capacité d'analyse de l'IHEDN et rendre plus lisible la contribution possible de ses associations face aux crises majeures.

\clearpage

\subsection{Organisation des acteurs : une gouvernance en triptyque}

Un tel projet nécessite des responsabilités clairement identifiées, afin de garantir un suivi méthodologique rigoureux.

Pour cela, une gouvernance en triptyque est proposée. Celle-ci se composerait d'abord d'un comité de pilotage, chargé d'assurer l'orientation stratégique, la validation des grandes étapes et le lien avec les instances décisionnelles. Un comité scientifique apporterait ensuite une validation méthodologique et contribuerait à la conformité juridique du dispositif. Enfin, la conception et la mise en \oe{}uvre seraient assurées par une équipe projet composée de membres du comité.

Une telle gouvernance, illustrée dans la figure \ref{fig:gouvernance-triptyque}, permettrait d'articuler expertise, légitimité institutionnelle et capacité de suivi, trois conditions nécessaires à la réussite d'une enquête de cette ampleur.

\begin{figure}[H]
\centering
\scalebox{1.5}{%
\begin{tikzpicture}[
    governance circle/.style={draw=black, line width=1.1pt, fill=gray!3},
    governance title/.style={font=\bfseries\footnotesize, align=center, text width=3.8cm},
    governance body/.style={font=\scriptsize, align=center, text width=3.6cm, anchor=north, inner sep=0pt},
    interaction arrow/.style={-{Stealth[length=2.4mm]}, line width=0.65pt},
    interaction label/.style={font=\tiny, align=center, fill=white, inner sep=1.2pt}
]
\def\governanceradius{2.05}
\coordinate (pilotage) at (0,2.15);
\coordinate (scientifique) at (-3.15,-1.25);
\coordinate (projet) at (3.15,-1.25);

\draw[governance circle] (pilotage) circle[radius=\governanceradius cm];
\draw[governance circle] (scientifique) circle[radius=\governanceradius cm];
\draw[governance circle] (projet) circle[radius=\governanceradius cm];

\draw[interaction arrow] (0.72,0.40) to[bend left=12]
    node[interaction label, pos=0.50, sloped, above] {Orientation\\Validation} (1.52,-0.05);
\draw[interaction arrow] (1.06,0.05) to[bend left=12]
    node[interaction label, pos=0.50, sloped, below] {Reporting} (0.26,0.50);

\draw[interaction arrow] (-0.72,0.40) to[bend right=12]
    node[interaction label, pos=0.50, sloped, above] {Cadrage\\Arbitrage} (-1.52,-0.05);
\draw[interaction arrow] (-1.06,0.05) to[bend right=12]
    node[interaction label, pos=0.50, sloped, below] {Avis\\Conformité} (-0.26,0.50);

\draw[interaction arrow] (-1.25,-1.05) to[bend left=10]
    node[interaction label, pos=0.50, above] {Appui méthodologique\\et RGPD} (1.25,-1.05);
\draw[interaction arrow] (1.25,-1.45) to[bend left=10]
    node[interaction label, pos=0.50, below] {Besoins terrain\\Retours opérationnels} (-1.25,-1.45);

\node[governance title] at (0,2.80) {Comité de pilotage};
\node[governance body] at (0,2.42) {%
    Présidents des trois\\régions pilotes\\[0.08cm]
    Présidence de l'Union IHEDN\\[0.08cm]
    Membres du comité d'étude
};

\node[governance title] at (-3.15,-0.60) {Comité scientifique};
\node[governance body] at (-3.15,-0.98) {%
    Expert en études\\et enquêtes sociologiques\\[0.08cm]
    Juriste RGPD
};

\node[governance title] at (3.15,-0.60) {Équipe projet};
\node[governance body] at (3.15,-0.98) {%
    Membres opérationnels\\du comité
};
\end{tikzpicture}
}
\caption{Organisation proposée de la gouvernance du projet de questionnaire.}
\label{fig:gouvernance-triptyque}
\end{figure}

Le comité a également identifié les principaux risques susceptibles d'affecter l'aboutissement du projet.

\clearpage

\subsection{Analyse des risques et facteurs de vigilance}

Le tableau \ref{tab:risques} présente les risques identifiés, en cohérence avec les retours d'expérience issus des travaux du comité.

\begin{table}[h]
\centering
\small
\setlength{\tabcolsep}{4pt}
\begin{tabular}{@{}>{\raggedright\arraybackslash}p{2.6cm}>{\raggedright\arraybackslash}p{7.2cm}C{2.1cm}C{1.7cm}@{}}
\toprule
\textbf{Risque} & \textbf{Description} & \textbf{Probabilité} & \textbf{Impact} \\
\midrule
Coordination & Difficulté de coordination entre les trois \AssociationsRegionalesIHEDN{} pilotes & Moyenne & Élevé \\
\addlinespace
Interprétation & Mauvaise lecture ou instrumentalisation des résultats, pouvant générer inertie ou incompréhension & Faible & Élevé \\
\addlinespace
Participation & Retards ou faibles taux de réponse, compromettant la représentativité des données & Moyenne & Élevé \\
\addlinespace
Périmètre & Confusion possible entre le potentiel des adhérents répondants, celui des associations et celui de l'ensemble des anciens auditeurs recensés & Moyenne & Élevé \\
\addlinespace
Ressources & Inadéquation potentielle entre les compétences requises pour assurer la qualité de conception et l'appropriation du questionnaire, et les ressources effectivement mobilisées & Faible & Élevé \\
\bottomrule
\end{tabular}
\caption{Cartographie des risques du projet.}
\label{tab:risques}
\end{table}

À ces risques s'ajoutent des points de vigilance déjà identifiés par le comité, notamment la disponibilité réelle des membres, la nécessité d'un socle commun de compétences, l'hétérogénéité des effectifs entre \AssociationsRegionalesIHEDN{} et l'importance d'une communication claire et pédagogique.

\begin{proposition}{Mettre en œuvre une démarche progressive et encadrée}
Le comité considère que la mise en place du questionnaire doit s'inscrire dans une démarche progressive, fondée sur une coordination étroite entre les associations pilotes, une communication claire auprès des membres et une vigilance particulière quant à la qualité méthodologique et à l'interprétation des résultats.
\end{proposition}

\clearpage

\subsection{Facteurs clefs de succès}

Les travaux du comité conduisent à identifier plusieurs conditions nécessaires à la réussite du projet. Celui-ci suppose d'abord une validation institutionnelle claire et un portage durable par les instances associatives concernées. Il repose également sur une conception méthodologique suffisamment rigoureuse pour garantir la crédibilité et l'exploitabilité des résultats recueillis.

L'adhésion des membres constitue un second enjeu majeur. Elle dépendra en partie de la capacité de l'association à expliquer les finalités du questionnaire, les modalités d'utilisation des données collectées et les bénéfices attendus pour les auditeurs comme pour les \AssociationsUnionIHEDN. L'analyse de l'annuaire 2022 rappelle que l'adhésion associative ne va pas de soi, y compris parmi les sessions récentes. La formulation du questionnaire devra donc permettre de comprendre les conditions concrètes de l'engagement, et non seulement de recueillir une disponibilité de principe.

Enfin, l'intérêt du dispositif reposera sur la capacité à produire une exploitation concrète des résultats, permettant de mieux identifier les compétences disponibles et de renforcer la lisibilité de l'\ARSeize{} auprès des autorités et partenaires institutionnels.

\begin{proposition}{Inscrire le questionnaire dans une démarche durable et exploitable}
Le comité considère que la réussite du questionnaire dépendra autant de sa crédibilité méthodologique que de sa capacité à produire des résultats concrètement exploitables pour les \AssociationsUnionIHEDN{} et leurs partenaires institutionnels.
\end{proposition}

Sur la base des éléments précédents, le comité a également étudié un calendrier prévisionnel en vue d'engager une telle démarche.

\clearpage

\subsection{Calendrier prévisionnel de mise en \oe{}uvre du questionnaire}

La figure \ref{fig:gantt-questionnaire-ar16} présente l'enchaînement des différentes étapes du projet, depuis leur planification jusqu'à l'obtention et l'exploitation des résultats du questionnaire en septembre 2027.

\begin{figure}[H]
\centering
\makeatletter
\providecommand{\ganttvrulelabel}[4][]{%
  \begingroup
  \ganttset{#1}%
  \gtt@tsstojulian{#4}{\gtt@vrule@slot}%
  \gtt@juliantotimeslot{\gtt@vrule@slot}{\gtt@vrule@slot}%
  \pgfmathsetmacro\y@upper{%
    \gtt@lasttitleline * \ganttvalueof{y unit title}%
  }%
  \pgfmathsetmacro\y@lower{%
    \gtt@lasttitleline * \ganttvalueof{y unit title}%
    + (\gtt@currentline - \gtt@lasttitleline - 1)%
    * \ganttvalueof{y unit chart}%
  }%
  \pgfmathsetmacro\x@mid{%
    (\gtt@vrule@slot - 1 + \ganttvalueof{vrule offset})%
    * \ganttvalueof{x unit}%
  }%
  \pgfmathsetmacro\y@label{\y@lower - #2}%
  \draw [/pgfgantt/vrule]
    (\x@mid pt, \y@upper pt) -- (\x@mid pt, \y@label pt)
    node [/pgfgantt/vrule label node] {#3};%
  \endgroup
}
\makeatother

\begin{ganttchart}[
    hgrid,
    x unit=0.0147cm,
    y unit chart=0.58cm,
    time slot format=isodate,
    title/.style={fill=gray!20, draw=gray!60},
    title label font=\bfseries\fontsize{6.4pt}{7pt}\selectfont,
    bar label font=\fontsize{6.8pt}{7.4pt}\selectfont,
    group label font=\bfseries\fontsize{6.8pt}{7.4pt}\selectfont,
    bar height=0.60,
    group height=0.38,
    group right shift=0,
    group left shift=0,
    group peaks tip position=0,
    group peaks height=0.18,
    group peaks width=0.2,
    bar/.style={fill=gray!35, draw=gray!70},
    group/.style={fill=gray!70, draw=gray!80},
    vrule/.style={dashed, line width=0.6pt, draw=black!70},
    vrule offset=.5,
    vrule label font=\fontsize{5pt}{5.6pt}\selectfont,
    vrule label node/.append style={anchor=north, rotate=45, align=left, inner sep=0.5pt, text=black!80},
    canvas/.style={draw=gray!50},
    y unit title=0.58cm,
    title height=1,
    title label anchor/.style={below=-1.6ex},
    bar label node/.append style={text width=4.2cm, align=right},
    group label node/.append style={text width=4.2cm, align=right}
]{2025-10-01}{2027-09-30}

\gantttitle{2025}{92}
\gantttitle{2026}{365}
\gantttitle{2027}{273} \\
\gantttitle{Oct}{31}
\gantttitle{Nov}{30}
\gantttitle{Déc}{31}
\gantttitle{Jan}{31}
\gantttitle{Fév}{28}
\gantttitle{Mar}{31}
\gantttitle{Avr}{30}
\gantttitle{Mai}{31}
\gantttitle{Juin}{30}
\gantttitle{Juil}{31}
\gantttitle{Août}{31}
\gantttitle{Sep}{30}
\gantttitle{Oct}{31}
\gantttitle{Nov}{30}
\gantttitle{Déc}{31}
\gantttitle{Jan}{31}
\gantttitle{Fév}{28}
\gantttitle{Mar}{31}
\gantttitle{Avr}{30}
\gantttitle{Mai}{31}
\gantttitle{Juin}{30}
\gantttitle{Juil}{31}
\gantttitle{Août}{31}
\gantttitle{Sep}{30} \\

\ganttgroup{Cadrage stratégique}{2025-10-01}{2026-06-25} \\
\ganttbar[bar/.append style={fill=blue!35}]{Étude bibliographique}{2025-10-01}{2026-04-30} \\
\ganttbar[bar/.append style={fill=blue!25}]{Sélection des trois AR pilotes}{2026-04-01}{2026-05-31} \\
\ganttbar[bar/.append style={fill=blue!45}]{Validation institutionnelle}{2026-03-30}{2026-05-25} \\

\ganttgroup{Conception du questionnaire}{2026-02-01}{2026-12-31} \\
\ganttbar[bar/.append style={fill=green!35}]{Objectifs, problématique et thèmes}{2026-02-01}{2026-12-31} \\
\ganttbar[bar/.append style={fill=green!25}]{Trame et grille d'exploitation}{2026-06-01}{2026-12-31} \\
\ganttbar[bar/.append style={fill=green!45}]{Maquette, support et outils}{2026-06-01}{2026-08-31} \\

\ganttgroup{Mobilisation et gouvernance}{2026-07-01}{2026-12-31} \\
\ganttbar[bar/.append style={fill=orange!35}]{Comité de pilotage}{2026-07-01}{2026-09-30} \\
\ganttbar[bar/.append style={fill=orange!25}]{Comité scientifique}{2026-07-01}{2026-09-30} \\
\ganttbar[bar/.append style={fill=orange!45}]{Ressources et partenariats}{2026-07-01}{2026-09-30} \\
\ganttbar[bar/.append style={fill=green!55}]{Validation RGPD et consentement}{2026-10-01}{2026-12-31} \\
\ganttbar[bar/.append style={fill=orange!55}]{Communication projet}{2026-10-01}{2026-12-31} \\
\ganttbar[bar/.append style={fill=orange!65}]{Préparation de l'emailing}{2026-10-01}{2026-12-31} \\

\ganttgroup{Lancement et déploiement}{2027-01-04}{2027-03-31} \\
\ganttbar[bar/.append style={fill=red!30}]{Lancement du questionnaire}{2027-01-04}{2027-02-28} \\
\ganttbar[bar/.append style={fill=red!40}]{Administration et relances}{2027-02-28}{2027-03-31} \\

\ganttgroup{Exploitation et valorisation}{2027-04-01}{2027-09-15} \\
\ganttbar[bar/.append style={fill=purple!25}]{Traitement des résultats}{2027-04-01}{2027-05-31} \\
\ganttbar[bar/.append style={fill=purple!35}]{Analyse et cartographie des compétences}{2027-04-01}{2027-06-30} \\
\ganttbar[bar/.append style={fill=purple!45}]{Rédaction du rapport de travail}{2027-06-01}{2027-08-31} \\
\ganttbar[bar/.append style={fill=purple!55}]{Capitalisation et restitution}{2027-08-01}{2027-09-15}

\ganttvrulelabel[vrule/.append style={draw=blue!70!black}, vrule label node/.append style={anchor=north east}]{10}{\shortstack[l]{15/05/2026\\Validation CODIR / Bureau}}{2026-05-15}
\ganttvrulelabel[vrule/.append style={draw=blue!70!black}, vrule label node/.append style={anchor=north west}]{34}{\shortstack[l]{25/06/2026\\Restitution du rapport de comité}}{2026-06-25}
\ganttvrulelabel[vrule/.append style={draw=red!70!black}, vrule label node/.append style={anchor=north west}]{34}{\shortstack[l]{31/03/2027\\Clôture de l'enquête}}{2027-03-31}
\ganttvrulelabel[vrule/.append style={draw=purple!70!black}, vrule label node/.append style={anchor=north east}]{10}{\shortstack[l]{15/09/2027\\Restitution finale}}{2027-09-15}

\end{ganttchart}

\caption{Calendrier prévisionnel de mise en \oe{}uvre du questionnaire de cartographie des compétences des membres de l'\ARSeize.}
\label{fig:gantt-questionnaire-ar16}
\end{figure}